% ===================================================================
%  تابلوی شمارهٔ ۱۲ — ایستگاه. --- راه‌آهن. --- کشتی‌ها.
%
%  Traduction, faite en 2026, en persan d'aujourd'hui, sur l'ido de
%  texto/io. Regles de la colonne : voir 10-tablo-01.tex. Une seule
%  numerotation.
%
%  LE SEUL « \textsc{} » DU LIVRET HORS DIALOGUE EST DANS CE TABLEAU,
%  ET C'EST LUI QUI A FIXE LE SEUIL DE kolonoj.py POUR LA COLONNE
%  JAVANAISE. L'ido y compose « \textsc{Noto.} --- » : une attribution
%  de forme, pas de locuteur. Le tableau 5 en compte trente-six,
%  celui-ci une, les quatorze autres aucune. La colonne persane n'a
%  pas besoin de ce seuil — ses dix regles lisent l'octet —, mais le
%  fichier le porte quand meme, et il faut le composer juste.
%
%  LE CHEMIN DE FER EST ARRIVE EN IRAN TRENTE ANS APRES JAVA, ET LE
%  VOCABULAIRE LE DIT. Le javanais du tableau 12 n'avait pas un mot a
%  lui : sepur, karcis, lokomotip, rel, wesel, peron, kondhektur,
%  masinis — tout neerlandais, parce que la premiere ligne de Java
%  date de 1867. Le persan a construit son transiranien en 1938, un
%  siecle apres l'invention, et il a eu le temps de composer :
%  راه‌آهن, « le chemin de fer », mot pour mot ; ایستگاه, « le lieu ou
%  l'on s'arrete » ; قطار, la file ; واگن et لکوموتیو seuls sont
%  empruntes, et au francais.
%
%  ET C'EST LA PREMIERE FOIS DU RELEVE QU'UNE DATE D'ARRIVEE SE LIT
%  DANS LA FORME DES MOTS. Ce qui arrive tot arrive avec son nom
%  etranger, parce qu'il n'y a rien pour le nommer ; ce qui arrive
%  tard trouve une langue prete a composer. Java a recu la machine
%  avant d'avoir un mot ; l'Iran l'a recue apres en avoir fait un. Le
%  meme objet mesure ici l'ecart de deux histoires.
%
%  UNE FAUTE D'ORDRE, LA SIXIEME DE LA COLONNE ET TOUJOURS LA MEME
%  ESPECE. Au bloc t12-06-1 la valise (24) avait ete posee avant la
%  sortie (32) vers laquelle on court. Le persan pose ordinairement
%  l'instrument avant le but, et c'est ce qui a fait glisser la main ;
%  la correction rejette « چمدان در دست » en fin de phrase, ou le
%  persan l'admet tres bien comme apposition. Six fautes en douze
%  tableaux, toutes de placement, aucune de grammaire : cette colonne
%  a un seul defaut et il est constant.
%
%  ET « ایستگاه » EST LE MOT LE PLUS PRODUCTIF DE LA COLONNE : il a
%  servi a la gare, puis a l'arret de tram (75) du tableau 10, puis a
%  la station thermale (15) du tableau 9, et il sert aujourd'hui a la
%  station-service et a la station spatiale. Un suffixe persan —
%  گاه, le lieu — attache a un verbe persan — ایستادن, s'arreter —
%  couvre tout ce que le francais partage entre gare, station, arret
%  et halte.
%
%  LE VOYAGE EST EN ALLEMAGNE, AVEC UNE FRONTIERE D'AVANT 1914 QUE
%  L'IDO SIGNALE LUI-MEME EN NOTE, ET LA COLONNE LA GARDE TELLE
%  QUELLE : on modernise la langue, jamais les choses, et une
%  frontiere est une chose.
% ===================================================================

%%K t12-apar-1 apar
\VUsaut{4.0mm}
\VUtitre{15.0pt}{60}{تابلوی شمارهٔ ۱۲}
\VUsaut{2.4mm}
\VUpk{11.4pt}{40}{ایستگاه. --- راه‌آهن. --- کشتی‌ها.}
\VUsaut{3.4mm}
\textsc{یادداشت.} --- \textit{جدول‌های ساعت که در هر کشور به کار
می‌روند گوناگون‌اند ؛ پس ساعت‌های} \VUgras{جدول فرانسوی}
\textit{را نمونه می‌گیریم.}

%%K t12-01-1 p
1. --- همین حالا از آلمان سفر کردم. پیش از رفتن، \VUgras{کتاب ساعت}
\textsuperscript{(15)} خریدم تا ساعت حرکت قطارها را بدانم. پس از آنکه
دریافتم مناسب‌ترین قطار ساعت نُه شب از پاریس حرکت می‌کند و ساعت یک و
سیزده دقیقه به استراسبورگ می‌رسد، سفر خود را آماده کردم، و فردای آن
روز مرا به ایستگاه بردند.

%%K t12-02-1 p
2. --- هنگامی که به سالن بزرگ حرکت درآمدم، پشت سر \VUgras{کشیش}
\textsuperscript{(2)} روستایی پیری که \VUgras{کیف سفر}
\textsuperscript{(3)} خود را به دست داشت، \VUgras{مسافران}
\textsuperscript{(1)} بسیاری از پیش رسیده بودند.

%%K t12-02-2 p
چون می‌خواستم \VUgras{تلگرافی} \textsuperscript{(5)} بفرستم، از
\VUgras{بازرس} \textsuperscript{(6)} خواستم \VUgras{ادارهٔ تلگراف}
\textsuperscript{(4)} را به من نشان دهد.

%%K t12-03-1 p
3. --- پیش از رفتن به \VUgras{سالن انتظار} \textsuperscript{(7)}، رفتم
تا \VUgras{بلیت} \textsuperscript{(9)} رفت و برگشت بگیرم.

%%K t12-04-1 p
4. --- کنار \VUgras{باجه} \textsuperscript{(8)} دیر نماندم، زیرا کسان
اندکی آنجا بودند. \VUgras{سربازی} \textsuperscript{(10)} که به مرخصی
می‌رفت، مهربانانه نوبت خود را به من داد، چنانکه وقت بس داشتم تا از
\VUgras{رئیس ایستگاه} \textsuperscript{(13)} چند آگاهی بخواهم که با
\VUgras{کمیسر} \textsuperscript{(12)} نگهبانی و با \VUgras{ژاندارم}
\textsuperscript{(11)} کشیک گفت‌وگو می‌کرد.

%%K t12-tit-1 sub
\VUsaut{3.0mm}
\VUpk{10.2pt}{40}{ثبت بار.}
\VUsaut{3.0mm}

%%K t12-05-1 p
5. --- پس از خریدن روزنامه در \VUgras{کتاب‌فروشی} \textsuperscript{(14)}،
\VUgras{بارهایم} \textsuperscript{(17)} را که \VUgras{باربر}
\textsuperscript{(16)} همین حالا با \VUgras{چرخ بار}
\textsuperscript{(19)} آورده بود وزن کردم ؛ \VUgras{کارمند}
\textsuperscript{(22)} گمارده بر \VUgras{ترازوی} \textsuperscript{(21)}
به من آگاهی داد که چمدانم سی و پنج کیلوگرم است — پس پنج کیلوگرم
\VUgras{اضافه بار} بود که باید در \VUgras{باجهٔ} بار
\textsuperscript{(23)} برایش پول افزوده بپردازم.

%%K t12-06-1 p
6. --- چون زود رسیده بودم، وقت آزاد داشتم تا رفت‌وآمد مسافران را در
ایستگاه بررسی کنم. برخی بارهای کم‌اهمیت خود را در \VUgras{امانت بار}
\textsuperscript{(25)} می‌گذاشتند یا پس می‌گرفتند ؛ دیگران
\VUgras{جدول‌های ساعت} \textsuperscript{(26)} را می‌خواندند. مسافران
تازه‌رسیده به سوی \VUgras{در خروج} \textsuperscript{(32)} می‌شتافتند،
\VUgras{چمدان} \textsuperscript{(24)} در دست. برخی در \VUgras{بخش
تحویل بار} \textsuperscript{(27)} در انتظار بودند و \VUgras{قبض}
\textsuperscript{(29)} خود را نشان می‌دادند تا چمدان‌ها،
\VUgras{صندوق‌ها} \textsuperscript{(28)} یا \VUgras{سبدهای}
\textsuperscript{(30)} خود را پس بگیرند.

%%K t12-07-1 p
7. --- \VUgras{مأموران عوارض} \textsuperscript{(33)} بسته‌ها را با دقت
بررسی می‌کردند تا ببینند آیا چیزی برای اظهار هست.

%%K t12-08-1 p
8. --- برخی مسافران به \VUgras{بوفه} \textsuperscript{(34)} رفتند،
جایی که \VUgras{پیشخدمت} \textsuperscript{(35)} با شور به آنان خدمت
می‌کرد. آنان باید بشتابند، زیرا \VUgras{قطار} \textsuperscript{(36)}
که سریع‌السیر است، دیر در ایستگاه نمی‌ماند.

%%K t12-09-1 p
9. --- سپس به سکوی حرکت رفتم و \VUgras{واگن} \textsuperscript{(37)}
مستقیمی جستم تا ناچار نشوم در میان راه واگن عوض کنم. به واگن راهرودار
بالا رفتم که \VUgras{کوپهٔ} \textsuperscript{(38)} سیگاری‌ها و کوپهٔ
ویژهٔ « بانوان تنها » دارد.

%%K t12-tit-2 sub
\VUsaut{3.0mm}
\VUpk{10.2pt}{40}{حرکت در شب.}
\VUsaut{3.0mm}

%%K t12-10-1 p
10. --- پس از بالا رفتن از \VUgras{رکاب} \textsuperscript{(40)}، بستن
\VUgras{در} \textsuperscript{(39)}، بالا زدن \VUgras{پرده}
\textsuperscript{(41)} و گذاشتن بارهای کوچکم روی \VUgras{تور}
\textsuperscript{(43)}، روی \VUgras{نیمکت} \textsuperscript{(42)}
نشستم. با دیدن \VUgras{چراغبان} \textsuperscript{(44)} که چراغ‌ها را
روشن می‌کرد، اندیشیدم که باید شبی را در واگن بگذرانم و مأمور گمارده
بر کرایهٔ \VUgras{بالش‌ها} \textsuperscript{(45)} را خواندم تا یکی
بخواهم.

%%K t12-11-1 p
11. --- مأمور قطار آمد تا بلیت‌ها را بررسی کند. از او پرسیدم چرا هنوز
حرکت نکرده‌ایم، زیرا ساعت درست است. پاسخ داد که \VUgras{رئیس قطار}
\textsuperscript{(46)} سرگرم گذاشتن دوچرخه‌ها در \VUgras{واگن بار}
\textsuperscript{(47)} است. وانگهی، هنوز همهٔ پاگرمکن‌ها
\textsuperscript{(48)} را عوض نکرده بودند.

%%K t12-12-1 p
12. --- اما به زودی حرکت می‌کردیم، زیرا \VUgras{آتشکار}
\textsuperscript{(50)} روی \VUgras{تندر} \textsuperscript{(49)} خود
زغال برداشته بود تا آتش \VUgras{لکوموتیو} \textsuperscript{(54)} را
برانگیزد، و \VUgras{ماشین‌چی} \textsuperscript{(51)} تنها در انتظار
علامت \VUgras{معاون رئیس ایستگاه} \textsuperscript{(52)} بود که بر
\VUgras{سکو} \textsuperscript{(53)} ایستاده بود.

%%K t12-13-1 p
13. --- \VUgras{دیگ بخار} \textsuperscript{(56)} لکوموتیو کر و کر
می‌کرد ؛ \VUgras{دودکش} \textsuperscript{(55)} سیل‌هایی از دود آمیخته
با \VUgras{بخار} \textsuperscript{(60)} بیرون می‌ریخت. \VUgras{سوت}
\textsuperscript{(59)} گوش‌خراشی به صدا درآمد و \VUgras{پیستون‌ها}
\textsuperscript{(58)} به جنبش درآمدند و \VUgras{چرخ‌های}
\textsuperscript{(57)} ماشین سنگین را راندند.

%%K t12-tit-3 sub
\VUsaut{3.0mm}
\VUpk{10.2pt}{40}{حرکت!}
\VUsaut{3.0mm}

%%K t12-14-1 p
14. --- شتاب قطار که بر \VUgras{ریل‌ها} \textsuperscript{(64)}
می‌دوید، تندتر می‌شد ؛ \VUgras{سکوها} \textsuperscript{(65)} ناپدید
شدند ؛ از \VUgras{مستراح‌ها} \textsuperscript{(66)} گذشتیم، و نیز از
رشته واگنی که بر \VUgras{خط} \textsuperscript{(63)} دیگر پارک شده بود :
\VUgras{واگن‌های خواب} \textsuperscript{(67)}، \VUgras{واگن رستوران}
\textsuperscript{(68)}، \VUgras{واگن پست} \textsuperscript{(69)}.

%%K t12-noto-1 noto
\textit{(*) یادداشت.} --- \textit{روشن است که شرح این سفر بر پایهٔ مرز
پیش از جنگ است.}

%%K t12-15-1 p
15. --- سپس، \VUgras{ایستگاه بار} \textsuperscript{(71)} از برابر
چشمان ما گذشت. زیر انبارها، مأموران \VUgras{کالاهایی}
\textsuperscript{(72)} را بار می‌کردند که با قطار کند فرستاده می‌شدند.
\VUgras{کارگران مانور} \textsuperscript{(76)} نزدیک \VUgras{مخزن آب}
\textsuperscript{(73)} \VUgras{واگن‌های دام} \textsuperscript{(75)} و
\VUgras{واگن‌های تخت} \textsuperscript{(79)} را مانور می‌دادند تا
\VUgras{قطار باری} \textsuperscript{(80)} را ببندند.

%%K t12-16-1 p
16. --- از واپسین \VUgras{دوراهی} \textsuperscript{(78)} گذشتیم که
سوزن‌بان نزدیک آن ایستاده بود ؛ از \VUgras{صفحهٔ علامت}
\textsuperscript{(86)} گذشتیم و ایستگاه در دور ناپدید شد.

%%K t12-17-1 p
17. --- به زودی از \VUgras{پل آهنی} \textsuperscript{(74)} گذشتیم که
از \VUgras{رود} \textsuperscript{(81)} می‌گذرد. پس از گذشتن از آن پل،
کنارهٔ رودخانه را دنبال کردیم، جایی که \VUgras{یدک‌کش‌های}
\textsuperscript{(82)} نیرومند \VUgras{کرجی‌های باری}
\textsuperscript{(83)} سنگین را می‌کشیدند. \VUgras{کشتی مسافربری}
\textsuperscript{(84)} را هم دیدیم، هنگامی که به \VUgras{اسکلهٔ
شناور} \textsuperscript{(85)} پهلو می‌گرفت.

%%K t12-18-1 p
18. --- پس از گذشتن بسیار تند از چند ایستگاه، از چند \VUgras{تونل}
\textsuperscript{(87)} گذشتیم و \VUgras{خانه‌های کوچک} بی‌شمار
\textsuperscript{(88)} \VUgras{نگهبانان راه‌بند} را پشت سر گذاشتیم.
تکان‌های قطار مرا می‌جنباند و به خواب ژرفی فرو رفتم.

%%K t12-tit-4 sub
\VUsaut{3.0mm}
\VUpk{10.2pt}{40}{ایستگاه دویچ‌آوریکور.}
\VUsaut{3.0mm}

%%K t12-19-1 p
19. --- ناگهان با آواز مأموری بیدار شدم که فریاد می‌زد :
« دویچ‌آوریکور! \textsuperscript{(*)}. همه پیاده شوید! » بازرسی گمرک و
همهٔ تشریفات ملال‌آور مرز پیش آمد. ناچار شدم ساعت جیبی خود را با
\VUgras{ساعت بزرگ} \textsuperscript{(89)} ایستگاه هماهنگ کنم که پنجاه
و پنج دقیقه جلوتر بود ؛ تازه بیدار شده بودم و به سختی \VUgras{عقربهٔ
بزرگ} \textsuperscript{(90)} را از \VUgras{کوچک}
\textsuperscript{(91)} بازمی‌شناختم ؛ خود \VUgras{صفحهٔ ساعت}
\textsuperscript{(92)} از مه بامدادی سراسر تار بود.

%%K t12-20-1 p
20. --- در حالی که مأموران گمرک بازرسی خود را می‌کردند، خمیازه‌کشان
\VUgras{اعلان‌هایی} \textsuperscript{(93)} را می‌خواندم که دیوارهای
ایستگاه را می‌آرایند. برای گذراندن وقت صبحانه خوردم، نقشهٔ
\VUgras{شبکهٔ} \textsuperscript{(94)} آلمان را نگاه کردم تا ببینم چه
اندازه راه هنوز مرا از استراسبورگ جدا می‌کند، و دوباره به واگن خود
درآمدم، نیرو گرفته از امید رسیدن زود به مقصد سفرم.

\VUsaut{6.0mm}
\VUfilet{22mm}
